% GENERATED FILE. Edit canonical lessons, then run npm run generate:books.
% canonical-source-hash: fnv1a64:18bc71cb38b93ef1
% canonical-lessons: TA-C441-aluku, TA-C441-naru, TA-C441-kuvi, TA-C441-marucularci-cey, TA-C441-uyirinam

\chapter{Rot and Recycling}
\label{ch:ta-rot-and-recycling}
\hlchaptermodality{441}

\begin{chapteropening}
\textbf{By the end of this chapter:} \emph{I can talk about things rotting, stinking and piling up, recycling, and living creatures.}

Complete the last lesson of chapter 441: I can talk about things rotting, stinking and piling up, recycling, and living creatures.
\end{chapteropening}

\section[\texorpdfstring{\ta{அழுகு} (aḻuku)}{அழுகு (aḻuku)}]{\ta{அழுகு} (aḻuku) — to rot}

\label{lesson:TA-C441-aluku}



\begin{quote}
Before the new one: say the Tamil for waste, then the Tamil for plastic.
\end{quote}

\subsection*{You'll want to know: \ta{அழுகு}}
\textbf{\ta{அழுகு}} — \emph{aḻuku} — \textquotedblleft{}to rot\textquotedblright{}.

\begin{quote}
\ta{சந்தையில்} \ta{தக்காளிகள்} \ta{அழுகின}.
\end{quote}

\begin{quote}
\emph{cantaiyil takkāḷikaḷ aḻukiṉa.}
\end{quote}

\begin{quote}
The tomatoes rotted in the market.
\end{quote}

Past (it …): \ta{அழுகியது}, \emph{aḻukiyatu}.

\ta{அழுகிய} \ta{பழம்} is a rotten fruit.

\begin{cousinweb}
\textbf{\ta{அழுகு}} (\emph{aḻuku}) is native Dravidian (DEDR 284), \textquotedblleft{}to rot, to decompose\textquotedblright{}.
\end{cousinweb}

\begin{grammarlens}[title={What you've built}]
One more word for getting around.
\end{grammarlens}

\subsection*{Your turn}
\begin{itemize}
\raggedright
  \item \emph{Say it:} \emph{aḻuku}
  \item \emph{Say it:} \emph{aḻuku}, once more
  \item \emph{Say it:} \emph{cantaiyil takkāḷikaḷ aḻukiṉa.} — The tomatoes rotted in the market.
  \item \emph{Recall:} say the Tamil for waste, then the Tamil for plastic, then say \emph{aḻuku} again
\end{itemize}

\begin{tcolorbox}[breakable,colback=teal!4,colframe=teal!35!black,title={Before you move on}]
What does \ta{அழுகு} mean? (\textquotedblleft{}To rot\textquotedblright{}.) Say it once more.
\end{tcolorbox}
\section[\texorpdfstring{\ta{நாறு} (nāṟu)}{நாறு (nāṟu)}]{\ta{நாறு} (nāṟu) — to stink}

\label{lesson:TA-C441-naru}



\begin{quote}
Before the new one: say the Tamil for plastic, then the Tamil for to rot.
\end{quote}

\subsection*{You'll want to know: \ta{நாறு}}
\textbf{\ta{நாறு}} — \emph{nāṟu} — \textquotedblleft{}to stink\textquotedblright{}.

\begin{quote}
\ta{குப்பைத்} \ta{தொட்டி} \ta{மோசமாக} \ta{நாறியது}.
\end{quote}

\begin{quote}
\emph{kuppait toṭṭi mōcamāka nāṟiyatu.}
\end{quote}

\begin{quote}
The rubbish bin stank badly.
\end{quote}

Past (it …): \ta{நாறியது}, \emph{nāṟiyatu}.

In old Tamil the verb also meant to smell sweet.

\begin{cousinweb}
\textbf{\ta{நாறு}} (\emph{nāṟu}) is native Dravidian (DEDR 2918), \textquotedblleft{}to give forth a smell, to stink\textquotedblright{}.
\end{cousinweb}

\begin{grammarlens}[title={What you've built}]
One more word for getting around.
\end{grammarlens}

\subsection*{Your turn}
\begin{itemize}
\raggedright
  \item \emph{Say it:} \emph{nāṟu}
  \item \emph{Say it:} \emph{nāṟu}, once more
  \item \emph{Say it:} \emph{kuppait toṭṭi mōcamāka nāṟiyatu.} — The rubbish bin stank badly.
  \item \emph{Recall:} say the Tamil for plastic, then the Tamil for to rot, then say \emph{nāṟu} again
\end{itemize}

\begin{tcolorbox}[breakable,colback=teal!4,colframe=teal!35!black,title={Before you move on}]
What does \ta{நாறு} mean? (\textquotedblleft{}To stink\textquotedblright{}.) Say it once more.
\end{tcolorbox}
\section[\texorpdfstring{\ta{குவி} (kuvi)}{குவி (kuvi)}]{\ta{குவி} (kuvi) — to pile up; to heap up}

\label{lesson:TA-C441-kuvi}



\begin{quote}
Before the new one: say the Tamil for to rot, then the Tamil for to stink.
\end{quote}

\subsection*{You'll want to know: \ta{குவி}}
\textbf{\ta{குவி}} — \emph{kuvi} — \textquotedblleft{}to pile up; to heap up\textquotedblright{}.

\begin{quote}
\ta{தெருவில்} \ta{குப்பை} \ta{குவிந்திருந்தது}.
\end{quote}

\begin{quote}
\emph{teruvil kuppai kuvintiruntatu.}
\end{quote}

\begin{quote}
Rubbish had piled up in the street.
\end{quote}

Past (it …): \ta{குவிந்தது}, \emph{kuvintatu}.

\ta{குவிந்தது} is it piled up; \ta{குவித்தேன்}, I heaped something up.

\begin{cousinweb}
\textbf{\ta{குவி}} (\emph{kuvi}) is native Dravidian (DEDR 1741), \textquotedblleft{}to crowd together, to be piled up\textquotedblright{}.
\end{cousinweb}

\begin{grammarlens}[title={What you've built}]
One more word for getting around.
\end{grammarlens}

\subsection*{Your turn}
\begin{itemize}
\raggedright
  \item \emph{Say it:} \emph{kuvi}
  \item \emph{Say it:} \emph{kuvi}, once more
  \item \emph{Say it:} \emph{teruvil kuppai kuvintiruntatu.} — Rubbish had piled up in the street.
  \item \emph{Recall:} say the Tamil for to rot, then the Tamil for to stink, then say \emph{kuvi} again
\end{itemize}

\begin{tcolorbox}[breakable,colback=teal!4,colframe=teal!35!black,title={Before you move on}]
What does \ta{குவி} mean? (\textquotedblleft{}To pile up; to heap up\textquotedblright{}.) Say it once more.
\end{tcolorbox}
\section[\texorpdfstring{\ta{மறுசுழற்சி} \ta{செய்} (maṟucuḻaṟci cey)}{மறுசுழற்சி செய் (maṟucuḻaṟci cey)}]{\ta{மறுசுழற்சி} \ta{செய்} (maṟucuḻaṟci cey) — to recycle}

\label{lesson:TA-C441-marucularci-cey}



\begin{quote}
Before the new one: say the Tamil for to stink, then the Tamil for to pile up.
\end{quote}

\subsection*{You'll want to know: \ta{மறுசுழற்சி} \ta{செய்}}
\textbf{\ta{மறுசுழற்சி} \ta{செய்}} — \emph{maṟucuḻaṟci cey} — \textquotedblleft{}to recycle\textquotedblright{}.

\begin{quote}
\ta{நாங்கள்} \ta{பழைய} \ta{காகிதத்தை} \ta{மறுசுழற்சி} \ta{செய்கிறோம்}.
\end{quote}

\begin{quote}
\emph{nāṅkaḷ paḻaiya kākitattai maṟucuḻaṟci ceykiṟōm.}
\end{quote}

\begin{quote}
We recycle old paper.
\end{quote}

Past (I …): \ta{மறுசுழற்சி} \ta{செய்தேன்}, \emph{maṟucuḻaṟci ceytēṉ}.

\ta{மறுசுழற்சி}, recycling, is the noun.

\begin{cousinweb}
\textbf{\ta{மறுசுழற்சி} \ta{செய்}} (\emph{maṟucuḻaṟci cey}) is built from \ta{மறு}, \textquotedblleft{}again\textquotedblright{}, and \ta{சுழற்சி}, \textquotedblleft{}a turning\textquotedblright{}, from \ta{சுழல்}, \textquotedblleft{}to spin\textquotedblright{} (DEDR 2698), with \ta{செய்}, \textquotedblleft{}to do\textquotedblright{}.
\end{cousinweb}

\begin{grammarlens}[title={What you've built}]
One more word for getting around.
\end{grammarlens}

\subsection*{Your turn}
\begin{itemize}
\raggedright
  \item \emph{Say it:} \emph{maṟucuḻaṟci cey}
  \item \emph{Say it:} \emph{maṟucuḻaṟci cey}, once more
  \item \emph{Say it:} \emph{nāṅkaḷ paḻaiya kākitattai maṟucuḻaṟci ceykiṟōm.} — We recycle old paper.
  \item \emph{Recall:} say the Tamil for to stink, then the Tamil for to pile up, then say \emph{maṟucuḻaṟci cey} again
\end{itemize}

\begin{tcolorbox}[breakable,colback=teal!4,colframe=teal!35!black,title={Before you move on}]
What does \ta{மறுசுழற்சி} \ta{செய்} mean? (\textquotedblleft{}To recycle\textquotedblright{}.) Say it once more.
\end{tcolorbox}
\section[\texorpdfstring{\ta{உயிரினம்} (uyiriṉam)}{உயிரினம் (uyiriṉam)}]{\ta{உயிரினம்} (uyiriṉam) — a living creature; an organism}

\label{lesson:TA-C441-uyirinam}



\begin{quote}
Before the new one: say the Tamil for to pile up, then the Tamil for to recycle.
\end{quote}

\subsection*{You'll want to know: \ta{உயிரினம்}}
\textbf{\ta{உயிரினம்}} — \emph{uyiriṉam} — \textquotedblleft{}a living creature; an organism\textquotedblright{}.

\begin{quote}
\ta{கடலில்} \ta{பல} \ta{வகை} \ta{உயிரினங்கள்} \ta{வாழ்கின்றன}.
\end{quote}

\begin{quote}
\emph{kaṭalil pala vakai uyiriṉaṅkaḷ vāḻkiṉṟaṉa.}
\end{quote}

\begin{quote}
Many kinds of creatures live in the sea.
\end{quote}

\ta{உயிர்}, life, + \ta{இனம்}, a kind.

\begin{cousinweb}
\textbf{\ta{உயிரினம்}} (\emph{uyiriṉam}) is built from \ta{உயிர்}, \textquotedblleft{}life\textquotedblright{} (DEDR 645), and \ta{இனம்}, \textquotedblleft{}a kind, a species\textquotedblright{} (DEDR 531).
\end{cousinweb}

\begin{grammarlens}[title={What you've built}]
One more word for getting around.
\end{grammarlens}

\subsection*{Your turn}
\begin{itemize}
\raggedright
  \item \emph{Say it:} \emph{uyiriṉam}
  \item \emph{Say it:} \emph{uyiriṉam}, once more
  \item \emph{Say it:} \emph{kaṭalil pala vakai uyiriṉaṅkaḷ vāḻkiṉṟaṉa.} — Many kinds of creatures live in the sea.
  \item \emph{Recall:} say the Tamil for to pile up, then the Tamil for to recycle, then say \emph{uyiriṉam} again
\end{itemize}

\begin{tcolorbox}[breakable,colback=teal!4,colframe=teal!35!black,title={Before you move on}]
What does \ta{உயிரினம்} mean? (\textquotedblleft{}A living creature; an organism\textquotedblright{}.) Say it once more.
\end{tcolorbox}
